\section{Method}
\label{sec:method}

\method has two parts. A world model $\WM$ predicts the DINOv2 tokens of an object at an arbitrary pose from a retrieved reference
view (\cref{sec:obs,sec:wm}); it is trained first and then frozen (\cref{sec:wmtrain}). A diffusion denoiser uses the prediction at
its current hypothesis at every step (\cref{sec:den}), and the same prediction ranks complete samples (\cref{sec:select}).
\Cref{fig:method} gives an overview.

\begin{figure}[t]
\centering
\figbox{5.2cm}{FIGURE PLACEHOLDER (architecture).\\[2pt]
(a) World model: reference tokens $\zref$ (+ geometry tokens, + 7 virtual-point tokens) $\rightarrow$ 6 adaLN transformer blocks
conditioned on Fourier features of $(g,\ g\,(\gref)^{-1})$ $\rightarrow$ predicted tokens $\zhat(g)$.\\
(b) Denoiser step: $g_t \rightarrow$ retrieval $\rightarrow \WM$ (frozen) $\rightarrow [\zobs,\zhat,\zhat-\zobs,w]$ tokens + conditioning
token $c(t,g_t,k)$ $\rightarrow$ 4 adaLN blocks $\rightarrow \epsilon_\theta$.\\
(c) Selection: $N$ samples, energy $E(g_i)$, $\argmin$.}
\caption{\textbf{Architecture of \method.} (a) The world model predicts tokens at pose $g$ from the nearest training view.
(b) Each denoising step compares this prediction at the current hypothesis with the observation. (c) The same prediction gives an
energy that ranks finished samples.}
\label{fig:method}
\end{figure}

\subsection{Observation tokens}
\label{sec:obs}
The input is a $224\times224$ crop $I$ of the object, cut and resized as in LiePose (dynamic zoom-in at training time), with
pixels outside the visible mask set to zero when a mask is part of the protocol (\cref{sec:protocol}). A frozen DINOv2
ViT-S/14~\cite{oquab2024dinov2} maps it to $16\times16$ patch tokens of dimension 384, which we average-pool $2\times2$ into
$\zobs\in\mathbb{R}^{64\times384}$. A per-token weight $w_p=w_0+m_p$ emphasises the object: $m_p\in\{0,1\}$ is the max-pooled visible mask
dilated by one token, and $w_0=0.05$ keeps a small weight on background tokens. Where no mask is available (SPEED+), $m_p=1$. The crop
is also described by box quantities $k$ (zoom, crop origin, box size and aspect), which contain no pose information. Tokens of all
training crops are computed once and cached; the image encoder is never trained.

\subsection{A retrieval-referenced world model}
\label{sec:wm}
Predicting the appearance of an object at an arbitrary pose from nothing is hard: the model would have to store the object's appearance
from every viewpoint. We instead predict from the nearest available view. For each object $o$, the \emph{reference bank} holds the
cached tokens, poses and crop boxes of the training crops, $\mathcal{B}_o=\{(\zobs_i,g_i,k_i)\}$, expanded by the object's symmetry
group $\mathcal{S}_o$ when one is known: each crop appears once per symmetric equivalent $(R_iS,t_i)$, $S\in\mathcal{S}_o$, subsampled to at
most 36 (T-LESS, YCB-V) or 60 (SYMSOL-T) equivalents. Given a query pose $g=(R,t)$, retrieval returns the nearest entry under the key
$[\mathrm{vec}(R_{\text{allo}}),\ \lambda t]$, $\mathrm{ret}(g)=(\zref,\gref,k^{\mathrm{ref}})$. While training both the world model and
the denoiser, entries that come from the observation's own source image are excluded, so the model never sees its target as a
reference. At test time the bank contains training crops only.

The world model predicts the tokens of the object at pose $g$, in the crop $k$ of the observation:
\begin{equation}
\zhat(g) = \WM\big(\zref,\ \gref \!\rightarrow\! g,\ k^{\mathrm{ref}},\ k\big)\in\mathbb{R}^{64\times384}.
\label{eq:wm}
\end{equation}
$\WM$ is a 6-block transformer of width 384 whose blocks are modulated (adaLN~\cite{peebles2023dit}) by Fourier features of the target
pose $g$ and of the relative pose $g\,(\gref)^{-1}$. Its input tokens are a linear projection of $\zref$.

\paragraph{Geometry tokens.}
A rotation of half a degree changes DINOv2 tokens very little, but it moves image points by a computable amount. When the camera
intrinsics and the crop box are available, we give this displacement to the model explicitly. For every reference token centre we
back-project onto a fronto-parallel plane at the reference object depth, move the point rigidly from $\gref$ to $g$, and project it into
the observation's crop; the token receives the displacement in token units and the log depth ratio, Fourier-encoded and added through a
zero-initialised layer. In addition, seven virtual object-frame points (the centre and six axis tips at a nominal radius; no CAD model)
are projected at both poses and prepended as tokens. We use geometry tokens on SYMSOL-T and SPEED+; on T-LESS and YCB-V the token cache
does not store the per-crop intrinsics they need, and the world model uses the reference tokens and pose only.
\TODO{state whether YCB-V uses geometry tokens once its world model is trained.}

\paragraph{Energy.}
The world model defines an energy of a pose hypothesis given the observation, the foreground-weighted squared token distance
\begin{equation}
E(g;\zobs) = \frac{\sum_p w_p\, \tfrac{1}{D}\|\zhat_p(g)-\zobs_p\|^2}{\sum_p w_p},
\label{eq:energy}
\end{equation}
with $D=384$ and the reference retrieved at $g$.

\subsection{Training the world model}
\label{sec:wmtrain}
We train $\WM$ so that its energy ranks poses, not only so that it reconstructs tokens. For each training crop with pose $g^\star$ we
form a candidate set $\mathcal{C}$: the ground truth, 11 uniformly random rotations, 12 rotations within $30^\circ$ of $g^\star$,
4 translation perturbations and 4 translations taken from other crops. The reference is retrieved at a noisy copy of $g^\star$
(rotation perturbed by $|\mathcal{N}(0,5^\circ)|$, translation by Gaussian noise), which mimics the poses at which the denoiser will
query it. With a soft target that accounts for symmetry,
$q_c \propto \exp\!\big(-d_{\mathcal{S}}(R_c,R^\star)^2/2\sigma_R^2-\|t_c-t^\star\|^2/2\sigma_t^2\big)$, where $d_{\mathcal{S}}$ is the rotation
distance to the nearest symmetric equivalent, the loss is
\begin{equation}
\mathcal{L}_{\mathrm{WM}} = \mathrm{CE}\Big(\mathrm{softmax}_c\big(-E(g_c)/\tau\big),\ q\Big) + \lambda_{\mathrm{rec}}\,E(g^\star),
\label{eq:wmloss}
\end{equation}
with a learned temperature $\tau$. The first term makes the energy rank candidates; the second anchors the prediction at the true
pose to the observed tokens. On SYMSOL-T we add regression targets from in-plane rotations of training images and from rendered pose
pairs (\cref{sec:supp-wm} of the supplement). The world model is trained once per dataset and is not updated afterwards.

\subsection{Comparing inside every denoising step}
\label{sec:den}
At step $t$, each sample with pose $g_t$ retrieves $\mathrm{ret}(g_t)$ and receives the world model's prediction $\zhat(g_t)$ in its
observation's crop. The denoiser forms one input token per image patch from the observation, the prediction, their difference and the
weight,
\begin{equation}
x_p=\mathrm{Lin}\big(\mathrm{LN}([\,\zobs_p,\ \zhat_p(g_t),\ \zhat_p(g_t)-\zobs_p,\ w_p\,])\big)+e_p,
\label{eq:tokens}
\end{equation}
with a learned position embedding $e_p$. A conditioning vector $c=\mathrm{Lin}(\mathrm{SiLU}(\gamma_t(t)+\gamma_g(r_t)+\gamma_k(k)))$ combines
Fourier embeddings $\gamma$ of the step, of the Lie-algebra coordinates $r_t$ of $g_t$, and of the crop box. $c$ is prepended as a token
and modulates four adaLN transformer blocks of width 256; the output at the conditioning token is mapped to the 6-dimensional noise
$\epsilon_\theta$. The prediction $\zhat$ is treated as a constant (no gradient reaches $\WM$).

LiePose's forward process, noise schedule, noise-prediction loss, sampler and weight averaging are unchanged; only $\epsilon_\theta$ is
replaced. Because $\zhat$ depends on $g_t$, the denoiser can read the error of its own hypothesis from the difference tokens instead of
inferring it from image features alone. Retrieval is a nearest-neighbour query on the CPU and the world model is a forward pass per
sample and step; we report the cost in \cref{sec:runtime}.

\paragraph{Control without the world model.}
To isolate the comparison, the control is the same network with $\zhat:=0$ in \cref{eq:tokens}, trained identically. It sees the same
DINOv2 observation tokens, crop box and pose embedding, so any difference between \method and the control is due to the predicted
tokens alone. Note that the control is a DINOv2 token denoiser, not LiePose, which uses a trained ResNet-34 encoder.

\subsection{Selecting among samples}
\label{sec:select}
A diffusion sampler returns one pose per run, and independent runs can land in different modes, \eg flipped configurations of a
near-symmetric object. We draw $N$ samples from independent priors and return $\hat g=\argmin_{i} E(g_i;\zobs)$
(\cref{eq:energy}), with the reference retrieved at each $g_i$. Selection uses only the observation and the training bank. $N$ is fixed
in advance for all benchmarks; results with one sample are the protocol of the main tables unless marked otherwise.
\TODO{fix N (e.g. 16) before running the selection experiments, and state it here.}

\subsection{Information used at test time}
\label{sec:protocol}
At test time \method uses the image crop, the box quantities $k$, the camera intrinsics (for geometry tokens), the object identity (to
select the reference bank $\mathcal{B}_o$, as in the BOP task where the target objects are given), and the cached tokens of training
crops. Crops are ground-truth boxes on all benchmarks: on T-LESS and YCB-V the visible box and visible mask, as in LiePose's
protocol~\cite{hsiao2024liepose}; on SPEED+ the box of the spacecraft keypoints projected with the ground-truth pose; on SYMSOL-T the
box of the rendered object's foreground. On T-LESS, YCB-V and SPEED+ every method we train uses the same crops. On SYMSOL-T LiePose
takes the full image, so the crop gives \method information that LiePose does not use; we report LiePose with the same crops as well
(\cref{tab:symsolt}). Rows whose crops come from a detector are marked. Apart from these crops, no ground-truth quantity enters
inference.
